% ==========================================================================
% Chapter 4 (expanded): Theoretical Model
% Governing equations for multiphase cavitating flow and the outlet-pressure
% derivation for a target cavitation number / supercavitation onset.
% Drop this into your existing Chapter 4 in place of / alongside Sections 4.1-4.3.
% ==========================================================================

\section{Governing equations of the multiphase cavitating flow}

The flow is treated as a homogeneous liquid-vapour mixture, consistent with the
\texttt{cavitatingFoam} barotropic-mixture solver used for the CFD case
(Chapter 5). The mixture density is
\begin{equation}
    \rho = \alpha_v \rho_v + \alpha_l \rho_l, \qquad \alpha_v + \alpha_l = 1,
\end{equation}
where $\alpha_v$ is the vapour volume fraction.

\paragraph{Mixture continuity.}
\begin{equation}
    \frac{\partial \rho}{\partial t} + \nabla\!\cdot(\rho \mathbf{U}) = 0
\end{equation}

\paragraph{Vapour volume-fraction transport.}
\begin{equation}
    \frac{\partial \alpha_v}{\partial t} + \nabla\!\cdot(\alpha_v \mathbf{U}) = \frac{\dot m}{\rho_v}
\end{equation}
where $\dot m$ is the interphase mass-transfer rate (evaporation/condensation).

\paragraph{Mixture momentum (RANS).}
\begin{equation}
    \frac{\partial (\rho \mathbf{U})}{\partial t} + \nabla\!\cdot(\rho \mathbf{U}\mathbf{U})
    = -\nabla p + \nabla\!\cdot\!\left[\mu_{\text{eff}}\left(\nabla \mathbf{U} + \nabla \mathbf{U}^T\right)\right]
\end{equation}
with $\mu_{\text{eff}} = \mu + \mu_t$, $\mu_t$ from the $k$-$\omega$ SST closure.

\paragraph{Barotropic equation of state} (linear compressibility model, as in
\texttt{thermodynamicProperties}):
\begin{align}
    \rho_l(p) &= \rho_{l,\text{sat}} + \psi_l (p - p_{\text{sat}}) \\
    \rho_v(p) &= \rho_{v,\text{sat}} + \psi_v (p - p_{\text{sat}})
\end{align}
Note $c_l = \psi_l^{-1/2} \approx 1508\ \text{m/s}$, matching the physical speed of
sound in water — confirming the compressibility constants are physically
grounded, not arbitrary numerical parameters.

\paragraph{Interphase mass transfer (Schnerr--Sauer, Rayleigh--Plesset based).}
\begin{equation}
    \dot m = C\,\frac{\rho_v \rho_l}{\rho}\,\alpha_v(1-\alpha_v)\,\frac{3}{R}
    \sqrt{\frac{2}{3}\frac{|p_v - p|}{\rho_l}}\ \operatorname{sign}(p_v - p)
\end{equation}
with bubble radius $R$ linked to $\alpha_v$ and the bubble number density $n$ by
\begin{equation}
    \alpha_v = \frac{n \tfrac{4}{3}\pi R^3}{1 + n \tfrac{4}{3}\pi R^3}.
\end{equation}
This source term is the mechanism that pins the local pressure to $p \approx p_v$
wherever a cavity forms; the cavitation number below follows directly from it
rather than being an independent assumption.

\section{Reduction to Bernoulli's equation outside the cavity}

Since $U_\infty = 25\ \text{m/s} \ll c_l$ (Mach $\approx 0.017$), the liquid region
is effectively incompressible and the momentum equation reduces along a
streamline to
\begin{equation}
    p + \tfrac12 \rho_l U^2 = p_\infty + \tfrac12 \rho_l U_\infty^2,
\end{equation}
giving the pressure coefficient
\begin{equation}
    C_p = \frac{p - p_\infty}{\tfrac12 \rho_l U_\infty^2} = 1 - \left(\frac{U}{U_\infty}\right)^2.
\end{equation}

\section{Cavity-surface condition and the cavitation number}

Setting $p = p_v$ (the value the mass-transfer source term enforces on the
cavity surface) in the pressure-coefficient relation:
\begin{equation}
    C_{p,\text{cavity}} = \frac{p_v - p_\infty}{\tfrac12 \rho_l U_\infty^2} = -\sigma,
    \qquad
    \sigma \equiv \frac{p_\infty - p_v}{\tfrac12 \rho_l U_\infty^2}.
\end{equation}

\section{Outlet pressure boundary condition and the target-$\sigma$ inversion}

The pressure-correction (PIMPLE) system requires exactly one Dirichlet
condition on $p$ to fix the absolute pressure level of the domain. The inlet
uses \texttt{zeroGradient} (Neumann) for $p$; the outlet, placed far downstream
where the flow has relaxed to the undisturbed ambient state, is the fixed-value
condition, so numerically $p_{\text{outlet}} \equiv p_\infty$.

Inverting the cavitation-number definition for a target $\sigma^*$:
\begin{equation}
    \boxed{p_{\text{outlet}} = p_v + \sigma^* \left(\tfrac12 \rho_l U_\infty^2\right)}
\end{equation}

For this case ($\rho_l = 1025\ \text{kg/m}^3$, $p_v = 3169\ \text{Pa}$,
$U_\infty = 25\ \text{m/s}$, $\tfrac12\rho U_\infty^2 = 320{,}312.5\ \text{Pa}$):

\begin{table}[h]
\centering
\begin{tabular}{ccl}
\toprule
$\sigma$ (target) & $p_{\text{outlet}}$ (Pa) & Regime \\
\midrule
0.20 & 67{,}200 & Partial cavity (staged start case) \\
0.10 & 35{,}200 & Developing cavity \\
0.05 & 19{,}200 & Target: supercavitation \\
\bottomrule
\end{tabular}
\caption{Outlet pressure required for each target cavitation number.}
\end{table}

\paragraph{Stability constraint.} The mixture sound speed follows Wood's
equation
\begin{equation}
    \frac{1}{\rho c^2} = \frac{\alpha_v}{\rho_v c_v^2} + \frac{\alpha_l}{\rho_l c_l^2},
\end{equation}
which collapses to only a few m/s in a two-phase cell. If $p_{\text{outlet}}
\le p_v$, the boundary cell itself has near-zero acoustic stiffness and the
fixed-value pressure BC becomes ill-posed, causing solver divergence. This is
why $\sigma$ must be stepped down in stages (0.20 $\to$ 0.10 $\to$ 0.05) from a
converged higher-$\sigma$ solution rather than initialized directly at the
target.

\section{Cavity size from a momentum/drag balance}

The cavitator drag coefficient follows the Reichardt closure
\begin{equation}
    C_d(\sigma) = C_{d0}(1+\sigma), \qquad C_{d0} \approx 0.8 \text{ for a sharp-edged disk}
\end{equation}
Since the pressure inside a developed cavity is essentially uniform at $p_v$,
the net pressure thrust over the cavity's maximum cross-section balances the
cavitator drag:
\begin{equation}
    D \approx \sigma \left(\tfrac12 \rho U_\infty^2\right) A_c
    = C_d(\sigma)\left(\tfrac12 \rho U_\infty^2\right) A_n
\end{equation}
which gives the maximum cavity diameter
\begin{equation}
    \boxed{D_c = d_n \sqrt{\frac{C_{d0}(1+\sigma)}{\sigma}}}
\end{equation}

For the $\varnothing 20\ \text{mm}$ disk cavitator ($d_n = 0.02\ \text{m}$)
against the $50\ \text{mm}$ body diameter:

\begin{table}[h]
\centering
\begin{tabular}{ccl}
\toprule
$\sigma$ & $D_c$ & vs.\ 50 mm body \\
\midrule
0.20 & 43.9 mm & cavity smaller than body \\
0.10 & 59.3 mm & cavity just exceeds body \\
0.05 & 82.0 mm & body fully enveloped -- supercavitation \\
\bottomrule
\end{tabular}
\caption{Predicted maximum cavity diameter vs.\ cavitation number.}
\end{table}

Cavity length does not reduce to a closed form from this simple momentum
balance; it requires the full slender-body cavity-shape solution (Garabedian,
1956). Validate $L_c$ either against the correlation in reference [3]
(Lee \& Kim, 2009, for disk/cone cavitators) or by measuring the
$\alpha_{\text{vapour}} = 0.5$ isosurface directly from the converged
$\sigma = 0.05$ CFD solution, per the Validation Plan (Chapter 7).
